% TOP_PROBLEM: {"id": 30005101, "title": "Erdős–Gallai Cycle-Decomposition Conjecture", "category_id": 3, "category": {"id": 3, "name": "graph_theory", "display_name": "Graph Theory", "description": "Problems involving graphs, networks, and their properties.", "slug": "graph-theory", "order_index": 3, "created_at": "2026-07-31T15:26:25.671Z"}, "status": "partially_solved"}
% FIRST_POSTED: 2026-10-06
% ATTEMPT_STATUS: solved
% OPENAI_CLAIM: {"schema_version":1,"source_repo":"https://github.com/openai/math","source_commit":"adc7f1241b42e322a6451854ab7e4b4c146bf78a","release_date":"2026-10-06","match":"full","resolution":"proved","independently_reviewed":false,"second_pass":"agreed","families":[{"family":"181","title":"The Erdős–Gallai cycle-decomposition conjecture","subject":"Combinatorics","manuscripts":[{"title":"A linear cycle-and-edge decomposition of every graph","pdf_path":"preprints/A-linear-cycle-and-edge-decomposition-of-every-graph-September-24-2026/main.pdf","readme_path":"preprints/A-linear-cycle-and-edge-decomposition-of-every-graph-September-24-2026/README.md","date":"2026-09-24","theorem_ref":"cor:eulerian-cycles (main.tex:117); thm:main (main.tex:62)","statement_sources":[{"path":"preprints/A-linear-cycle-and-edge-decomposition-of-every-graph-September-24-2026/build/main.tex","line":117,"label":"cor:eulerian-cycles"},{"path":"preprints/A-linear-cycle-and-edge-decomposition-of-every-graph-September-24-2026/build/main.tex","line":62,"label":"thm:main"}]}],"lean":{"doc_path":"lean/docs/181.md","comparators":["lean/ComparatorChallenges/CycleDecomposition.json","lean/ComparatorChallenges/CycleDecomposition.lean"],"declarations":["OAI.ErdosGallai.erdos_gallai"],"covers_main_theorem":true},"reasoning_trace":null}]}
\documentclass[11pt]{article}
\usepackage[margin=25mm]{geometry}
\usepackage{fontspec}
\setmainfont{Latin Modern Roman}
\usepackage{amsmath,amssymb,mathtools}
\usepackage{unicode-math}
\setmathfont{Latin Modern Math}
\usepackage{enumitem}
\usepackage{xurl,hyperref,bookmark}
\hypersetup{colorlinks=true,urlcolor=blue}
\setcounter{secnumdepth}{0}
\setlength{\parindent}{0pt}
\setlength{\parskip}{5pt}
\setlength{\emergencystretch}{3em}
\begin{document}
\section*{Erdős–Gallai Cycle-Decomposition Conjecture}

\subsection{OpenAI's claim}
\textit{Linear cycle-and-edge decomposition}: Every finite simple graph on $n$ vertices is claimed to have its edges partitioned using no more than $Cn$ pieces, each either a simple cycle or a single edge. The constant is the same for all graph orders.

\subsection{Mathematical statement of the claimed result}
\paragraph{Mathematical claim 1.} There is an absolute $C>0$ such that every finite simple undirected $n$-vertex graph admits an edge partition into at most $Cn$ pieces, each a single edge or the edge set of a cycle. For an Eulerian graph the corollary gives at most $Cn$ cycles alone. It further claims that, for every fixed $\delta,p>0$, some $\epsilon,n_0$ make the sharper \[\text{number of cycles}\le\Delta(G)/2+\delta n\] possible for Eulerian $G$ whenever $n\ge n_0$ and every bipartition $A\mathbin{\dot\cup}B=V(G)$ with $|A|,|B|\ge\epsilon n$ has $e_G(A,B)\ge p|A||B|$. The unconditional claim remains a linear bound with an unspecified constant. Eulerian means that every vertex has even degree; connectivity is unnecessary.

\subsection{Outline of the argument}
\textit{Linear cycle-and-edge decomposition}: Long cycles are removed across a sequence of decreasing scales and the residue is assigned to overlapping boxes with controlled total order. Expanding pieces are extracted within the boxes. A carefully chosen layer pays for a whole prefix of scales, while reserved routing edges and small residual graphs keep the total decomposition cost linear.

\subsection{Scope relative to this problem}
\textit{Linear cycle-and-edge decomposition}: The theorem partitions all edges of every finite simple n-vertex graph into at most C n cycles and single edges with one universal C, exactly the linear decomposition assertion.

With the hypotheses and conventions above, the claimed result addresses the whole mathematical assertion on this page.

\subsection{Formal verification}
Family 181 describes Lean coverage for a uniformly linear number of cycles and singleton edges in an edge partition.

\textit{Linear cycle-and-edge decomposition}: The manuscript links \texttt{CycleDecomposition.lean} (declaration \texttt{OAI.ErdosGallai.erdos\_gallai}).

The documented selection includes this manuscript’s main assertion; that assertion retains the restrictions stated above.

This site has not run the Lean development. These links describe OpenAI’s selected formal statements and should be read with their stated scope.

\subsection{Status caveat}
OpenAI reports that an internal model produced the manuscripts. Their arguments have not been independently verified as proofs by this site. Any portion outside the documented formal scope may contain mistakes, and the catalogue label reports the announced mathematical claim.

\subsection{Sources}
\begin{itemize}
\item OpenAI, \emph{A linear cycle-and-edge decomposition of every graph}, 2026-09-24; citation key \texttt{OAI:A-linear-cycle-and-edge-decomposition-of-every-graph-September-24-2026}. \href{https://github.com/openai/math/blob/adc7f1241b42e322a6451854ab7e4b4c146bf78a/preprints/A-linear-cycle-and-edge-decomposition-of-every-graph-September-24-2026/main.pdf}{PDF (GitHub)}; \href{https://github.com/openai/math/raw/adc7f1241b42e322a6451854ab7e4b4c146bf78a/preprints/A-linear-cycle-and-edge-decomposition-of-every-graph-September-24-2026/main.pdf}{PDF (direct)}; \href{https://github.com/openai/math/blob/adc7f1241b42e322a6451854ab7e4b4c146bf78a/preprints/A-linear-cycle-and-edge-decomposition-of-every-graph-September-24-2026/README.md}{Manuscript page and citation}; \href{https://github.com/openai/math/blob/main/preprints/A-linear-cycle-and-edge-decomposition-of-every-graph-September-24-2026/main.pdf}{Latest version}; Statement source: \href{https://github.com/openai/math/blob/adc7f1241b42e322a6451854ab7e4b4c146bf78a/preprints/A-linear-cycle-and-edge-decomposition-of-every-graph-September-24-2026/build/main.tex\#L117}{cor:eulerian-cycles}; Statement source: \href{https://github.com/openai/math/blob/adc7f1241b42e322a6451854ab7e4b4c146bf78a/preprints/A-linear-cycle-and-edge-decomposition-of-every-graph-September-24-2026/build/main.tex\#L62}{thm:main}.
\item Family 181: \href{https://github.com/openai/math/blob/adc7f1241b42e322a6451854ab7e4b4c146bf78a/lean/docs/181.md}{Lean scope}.
\item \href{https://github.com/openai/math/blob/adc7f1241b42e322a6451854ab7e4b4c146bf78a/lean/ComparatorChallenges/CycleDecomposition.json}{Comparator CycleDecomposition.json}.
\item \href{https://github.com/openai/math/blob/adc7f1241b42e322a6451854ab7e4b4c146bf78a/lean/ComparatorChallenges/CycleDecomposition.lean}{Comparator CycleDecomposition.lean}.
\item \href{https://github.com/openai/math/blob/adc7f1241b42e322a6451854ab7e4b4c146bf78a/CONTENTS.md}{OpenAI catalogue entry}.
\end{itemize}
\end{document}
